\chapter{Rehearsal}
\label{ch:rehearsal}

\section*{What this chapter is for}

Reading Part~\ref{part:build} does not prepare you for a build round. One timed
rehearsal with another person in the room does more than all of it, and this
chapter is the protocol.

\section{Why it cannot be done alone}

The things that go wrong are interpersonal and time-bound: asking too few
questions, not noticing a checkpoint, going silent under pressure, over-building
before anything runs, failing to notice that the person has been trying to tell
you something.

None is observable from the inside. \judgement{A solo run measures whether you
can build the thing, which is the one question the round is not asking.}

\section{The protocol}

\begin{kitmethod}
\textbf{Find a person.} They do not need to be an engineer. They need to be
able to answer questions in character and hold back what they have been told to
hold back.

\textbf{Give them a role brief in a domain you do not know.} Four are supplied
in \code{drills/role-briefs/}. Not your own domain \dash{} the whole point is
that you have to ask.

\textbf{Full length, timed, screen shared, recorded.} Not a shortened version.
The failure modes appear in the second half.

\textbf{Do not stop when something breaks.} That is the most valuable part of
the recording.

\textbf{Do not watch it back for half an hour.} You will not be able to see it
usefully before then.
\end{kitmethod}

\section{The role brief, and what makes it work}

Each brief has four parts, and the third is what makes the rehearsal worth
doing.

\begin{itemize}
  \item \textbf{What the owner knows} and will say when asked.
  \item \textbf{How they behave} \dash{} vague about volumes, specific about
        irritations, uninterested in technology.
  \item \textbf{What they do not say until asked outright.} There is always a
        second actor and a constraint they have not mentioned. If you do not
        ask, you do not find out, and you build the wrong thing. This is the
        instrument.
  \item \textbf{Four timed interjections} they hold and you do not know about:
        two changes of requirement and two challenges, at fixed points,
        including one after the build has started.
\end{itemize}

\judgement{The mid-build interjection is the most informative sixty seconds of
the whole rehearsal. What you do when a requirement changes at minute
sixty-five \dash{} whether you absorb it, negotiate it, or quietly ignore it
\dash{} is the thing a real session will find out about you, and most people
have never observed themselves doing it.}

\section{The scoring table}

Score afterwards, from the recording, not from memory.

\begin{kittable}{@{}lXl@{}}
\textbf{\#} & \textbf{Measure} & \textbf{Target} \\
\midrule
1 & first ``why'' question & by minute 3 \\
2 & questions asked before summarising & 6--9 \\
3 & discovery closed & by the checkpoint \\
4 & written scope artefact exists & yes \\
5 & second actor discovered & yes \\
6 & first line of real code & by the checkpoint \\
7 & silences longer than 20 seconds & 0 \\
8 & assistant outputs rejected aloud & $\geq 2$ \\
9 & the hard case was built and checked & yes \\
10 & demonstration in the product & yes \\
11 & said what you were unsure about & yes \\
12 & checkpoints passed without noticing & 0 \\
\end{kittable}

\begin{kitnote}
One row matters more than the others: \textbf{row 3}. When discovery closed
predicts most of the rest of the session, because everything downstream
inherits its slippage \dash{} and unlike the others it is a single number you
can read off the recording without judgement.

If you only track one thing across several rehearsals, track that.
\end{kitnote}

\section{What to do with the result}

Take the three worst rows. Drill those, twenty minutes each, in isolation
\dash{} not by running the whole thing again.

\begin{kittable}{@{}lX@{}}
\textbf{Weak row} & \textbf{The isolated drill} \\
\midrule
1, 2 & ten minutes of discovery only, against a new brief, no building \\
3 & the same, with a hard stop at the checkpoint \\
7 & narrate a build you already know how to do, alone, watching for silences \\
8 & thirty minutes of assisted building, counting rejections \\
9 & build one hard case, end to end, with its check, nothing else \\
10, 11 & the closing five minutes only, repeatedly \\
\end{kittable}

\judgement{Running the whole rehearsal again is the appealing option and the
less useful one. A second full run mostly re-measures; twenty minutes on the
weakest row changes it.}

\section{Rehearsing the other two formats}

\textbf{Take-home:} set a three-hour timer on an unfamiliar brief and keep to
it. Score only two things: was there a failing check in the first half-hour,
and did you stop on time.

\textbf{AI-free:} build something small without assistance, timed. If you have
been working with assistance daily, this is slower than you expect and the gap
is not predictable from the inside.

\begin{drill}
Book the rehearsal before you finish this chapter. It is the single
highest-value hour in the book and it is the one thing in it that requires
somebody else's diary, which is why it is the thing that does not happen.
\end{drill}

\section*{What this chapter does not claim}

The targets in the scoring table are the author's, set at values that a
well-run session hits rather than derived from observed sessions. Treat them as
a consistent yardstick across your own rehearsals rather than as a standard
somebody else is applying.

The claim that one rehearsal beats the rest of Part~\ref{part:build} is a
judgement. It rests on the observation that the failure modes are
interpersonal, and that reading about an interpersonal failure mode has never
been a reliable way to stop having it.
